% =============================================================================
% Problem Template
% =============================================================================
% Copy this file into:
%   latex/sections/<domain>/problems/lc-XXX-problem-name.tex
%
% Fill in the metadata below and remove any sections not required.
% =============================================================================

% STATUS: draft
% PROMOTE: exemplars
% TAGS: strings, two pointers
% DIFFICULTY: easy
% SOURCE: LeetCode 125

\ProblemTitle{Valid Palindrome}
\ProblemMeta{Easy}{Strings, Two Pointers}{LC-125}
\label{prob:lc-125-valid-palindrome}

\ProblemSummary
Given a string $s$, determine whether it is a palindrome when considering only alphanumeric characters and ignoring case.
Non-alphanumeric characters (punctuation/whitespace) are ignored. After filtering, the empty string is considered a palindrome.

\KeyIdea
We need to compare the string from both ends; a two-pointer approach lets us scan inward while skipping non-alphanumeric characters, achieving $O(n)$ time and $O(1)$ extra space.

% -----------------------------------------------------------------------------
% Optional: Author Thinking (excluded from exemplar builds)
% -----------------------------------------------------------------------------
\begin{Thinking}
<Informal reasoning, failed approaches, intuition-building, or mental models.
This section is intentionally exploratory and may be omitted in promoted or
exemplar versions of the book.>
\end{Thinking}

% -----------------------------------------------------------------------------
% Algorithm
% -----------------------------------------------------------------------------
\Algorithm
\begin{CodeBlock}[style=pseudo,caption={Pseudocode}]{14}
    initialise left = 0, right = length(s) - 1
    while left < right:
        if s[left] is not alphanumeric:
            left = left + 1
            continue
        if s[right] is not alphanumeric:
            right = right - 1
            continue
        if lowercase(s[left]) != lowercase(s[right]):
            return false
        left = left + 1
        right = right - 1
    return true
\end{CodeBlock}

% -----------------------------------------------------------------------------
% Correctness
% -----------------------------------------------------------------------------
\Correctness
At the start of each iteration, all alphanumeric characters strictly outside the interval $[left, right]$ have already been matched (case-insensitively) with their corresponding partner on the other side.
Skipping a non-alphanumeric character is safe because such characters are ignored by the problem definition.
If $\text{lower}(s[left]) \neq \text{lower}(s[right])$, then an unmatched pair exists, so the string cannot be a valid palindrome and we return \texttt{False}.
When the loop terminates ($left \ge right$), every relevant pair has been verified; a middle character (odd length after filtering) does not affect palindromicity.

% -----------------------------------------------------------------------------
% Complexity
% -----------------------------------------------------------------------------
\Complexity
Each pointer moves inward at most $n$ times, so time complexity is $O(n)$.
The algorithm uses $O(1)$ extra space (in-place scanning; no filtered string constructed).

% -----------------------------------------------------------------------------
% Implementation
% -----------------------------------------------------------------------------
%\bigskip
%\noindent\textbf{C++ Implementation}

%\lstinputlisting[
%  style=cpp,
%  caption={C++ Solution: <Problem Title>}
%]{/../problems/<domain>/<difficulty>/lc-XXX-problem-name/solution.cpp}

\bigskip
\noindent\textbf{Python Implementation}

\lstinputlisting[
  language=Python,
  backgroundcolor=\color{codebg},
  basicstyle=\ttfamily\small,
  frame=single,
  breaklines=true,
  caption={Python Solution: <Problem Title>}
]{../problems/<domain>/<difficulty>/lc-XXX-problem-name/solution.py}

\ProblemDivider
